\section{Elementary estimates}
\label{sec:perturbations_estimates}

\begin{lemma}[Exponential sums on windows]
    \label{lem:ldp_vandermonde_window}
    \lean{Complex.exists_window_le_norm_sum_mul_pow}
    \leanok
    Let $\mu_1,\dots,\mu_K$ be distinct complex numbers of modulus one and
    $a\in\C^K$ nonzero. There is $c_0>0$ such that for every integer $t\ge0$
    some $u\in\{0,\dots,K-1\}$ satisfies
    $\bigl\lvert\sum_{j=1}^Ka_j\mu_j^{\,t+u}\bigr\rvert\ge c_0$.
\end{lemma}
\begin{proof}
    \leanok
    The Vandermonde map $v\mapsto(\sum_jv_j\mu_j^{\,u})_{0\le u<K}$ is
    injective, so it is bounded below on $\C^K$. The vector
    $(a_j\mu_j^t)_j$ has the maximum norm of $a$.
\end{proof}

\begin{lemma}[Expectations in overlapping states]
    \label{lem:ldp_expectation_overlap}
    \lean{LinearMap.IsSymmetric.norm_inner_mul_norm_sub_le}
    \leanok
    For a Hermitian operator $Z$ and unit vectors $\phi,\psi$, with
    $\langle Z\rangle_\chi=\bra{\chi}Z\ket{\chi}$ and
    $\sigma_\chi(Z)=\lVert(Z-\langle Z\rangle_\chi)\chi\rVert$,
    \begin{align}
        \lvert\braket{\phi}{\psi}\rvert\,
        \lvert\langle Z\rangle_\phi-\langle Z\rangle_\psi\rvert
         & \le\sigma_\phi(Z)+\sigma_\psi(Z).
        \label{eq:ldp_expectation_overlap}
    \end{align}
\end{lemma}
\begin{proof}
    \leanok
    Since $\langle Z\rangle_\phi$ is real,
    $(\langle Z\rangle_\phi-\langle Z\rangle_\psi)\braket{\phi}{\psi}
        =\bra{\phi}(Z-\langle Z\rangle_\psi)\ket{\psi}
        -\bra{\phi}(Z-\langle Z\rangle_\phi)\ket{\psi}$, and applying the
    Cauchy--Schwarz inequality to both terms yields the bound.
\end{proof}

\begin{lemma}[Word values of a functional under two spanning families]
    \label{lem:block_symmetry_phase_power}
    \lean{LinearMap.map_list_prod_eq_pow_mul_of_length_three}
    \leanok
    Let $\phi$ be a linear functional on a complex algebra $\mathcal A$ with
    $\phi(1)\neq0$ such that $\phi(MN)=0$ for all $N$ forces $M=0$, and let
    $(x_i)_i$ and $(y_i)_i$ be finite families spanning $\mathcal A$. If
    $\phi(y_w)=c_n\phi(x_w)$ for every word $w$ of length $n\in\{1,2,3\}$,
    where $x_w$ and $y_w$ denote the ordered products along $w$, then
    $\phi(y_w)=c_1^{|w|}\phi(x_w)$ for every word $w$.
\end{lemma}
\begin{proof}
    \leanok
    Nondegeneracy against the spanning family $(y_i)$ makes $x_i\mapsto y_i$
    a well-defined linear map $L$: if $\sum_ia_ix_i=0$, then
    $\phi\bigl((\sum_ia_iy_i)y_j\bigr)=c_2\phi\bigl((\sum_ia_ix_i)x_j\bigr)=0$
    for every $j$. By multilinearity,
    $\phi(LX)=c_1\phi(X)$, $\phi(LX\,LY)=c_2\phi(XY)$ and
    $\phi(LX\,LY\,LZ)=c_3\phi(XYZ)$ for all $X,Y,Z$. The map $L$ is onto, so
    pairing against $LZ$ gives $LX\,LY=L(XY)L(1)$ and $c_1L(1)=c_2\cdot1$.
    Since $\phi(1)\neq0$ lies in the image of $\phi\circ L=c_1\phi$, we have
    $c_1\neq0$; applying $\phi$ to $c_1L(1)=c_2\cdot1$ gives $c_2=c_1^2$,
    hence $L(1)=c_1\cdot1$ and $LX\,LY=c_1L(XY)$. Induction on the word
    gives $y_w\,L(1)=c_1^{|w|}L(x_w)$, and applying $\phi$ proves the claim.
\end{proof}

\begin{lemma}[Enumerations listing a subset first]
    \label{lem:perturbations_enumeration_subset_first}
    \lean{Finset.exists_equiv_val_lt_card_iff_mem}
    \leanok
    Let $X$ be a finite set with $n$ elements and let $S\subseteq X$. There is
    a bijection $\pi:\{0,\dots,n-1\}\to X$ with $\pi(k)\in S$ if and only if
    $k<|S|$.
\end{lemma}
\begin{proof}
    \leanok
    Enumerate $S$ by $\{0,\dots,|S|-1\}$ and its complement by the remaining
    indices.
\end{proof}

\begin{lemma}[Rank bound for one isometric family]
    \label{thm:seqanc_rank_bound}
    \lean{Matrix.card_le_mul_of_sum_conjTranspose_mul_eq_one}
    \leanok
    If $A^0,\dots,A^{d-1}$ are $a\times b$ matrices with
    $\sum_i(A^i)^\dagger A^i=\Id_b$, then $b\le ad$.
\end{lemma}
\begin{proof}
    \leanok
    The $(ad)\times b$ matrix $S$ obtained by stacking the $A^i$ satisfies
    $S^\dagger S=\Id_b$, so $b=\operatorname{rank}(S^\dagger S)
        =\operatorname{rank}S\le ad$.
\end{proof}

\begin{lemma}[The radical of a nonnegative quadratic form]
    \label{lem:perturbations_quadratic_form_radical}
    \lean{QuadraticForm.polar_eq_zero_of_nonneg, QuadraticForm.mem_radical_iff_of_nonneg,
        QuadraticForm.radical_ne_top_of_ne_zero}
    \leanok
    Let $Q$ be a real quadratic form with $Q(x)\ge0$ for every $x$. A zero $x$
    of $Q$ is orthogonal to every vector for the polar form, so the zero set of
    $Q$ is its radical, a linear subspace. If $Q\neq0$ this subspace is proper.
\end{lemma}
\begin{proof}
    \leanok
    If $Q(x)=0$ and $p$ is the polar value of $x$ and $y$, then for
    $t=-p/(Q(y)+1)$ one has $Q(x+ty)=-p^2/(Q(y)+1)^2\ge0$, so $p=0$. The radical
    of the whole space would make $Q$ vanish identically.
\end{proof}

\begin{lemma}[Compactness of finite convex weights]
    \label{thm:finite_convex_weights_compactness}
    \lean{exists_convex_weights_subsequence,
        exists_convex_weights_of_approximate_evaluations}
    \leanok
    Let $J$ be finite, and let $w_j(n)$ be real weights that eventually lie
    in $[0,1]$, with $\sum_jw_j(n)\to1$. There are a strictly increasing
    subsequence $n_s$ and weights $p_j\geq0$, $\sum_jp_j=1$, such that
    $w_j(n_s)\to p_j$ for every $j$.
    For any family of tests $x$, if
    $f(x)-\sum_jw_j(n)g_j(x)\longrightarrow0$ for every test $x$,
    then one choice of normalized weights satisfies
    $f(x)=\sum_jp_jg_j(x)$ for all tests. The test family need not be countable.
\end{lemma}
\begin{proof}\leanok
    Extract a convergent subsequence in the compact cube $[0,1]^J$.
    Continuity of the finite sum shows that the limit has total mass one.
    For each test, compare the assumed limit with the evaluation along
    this same subsequence. Uniqueness of limits gives the asserted identity.
\end{proof}

\begin{lemma}[Uniform upper bounds from all selections]
    \label{thm:uniform_upper_bound_all_selections}
    \lean{Filter.eventually_forall_le_of_tendsto_all_selections}
    \leanok
    Let $X$ be a nonempty set, let $\mathcal F$ be a filter on an index set
    $I$, and let $e:I\times X\to\mathbb R$.
    Suppose $e(i,x_i)\to0$ along $\mathcal F$ for every function
    $i\mapsto x_i$ from $I$ to $X$. Then, for every $\varepsilon>0$,
    eventually $e(i,x)\leq\varepsilon$ simultaneously for every $x\in X$.
    No topology or countability assumption on $X$ is required, and $e$
    need not be nonnegative. This elementary implication applies to the
    uniform projection estimates in
    \cite[Lemma~commutation(ii)]{Nachtergaele1996Spectral}.
\end{lemma}
\begin{proof}\leanok
    For each $i$, select a parameter violating the upper bound if one exists;
    otherwise select any element of $X$.
    Convergence along this selection gives $e(i,x_i)<\varepsilon$ eventually.
    The choice of $x_i$ then implies $e(i,x)\leq\varepsilon$ for every $x$.
\end{proof}
